\subsection{Bipolar Junction Transistor (BJT)}
\markright{Bipolar Junction Transistor (BJT)}

\subsubsection*{Definition}
A bipolar junction transistor (BJT) is a three-terminal semiconductor device that
controls a large collector current $I_C$ with a small base current $I_B$ \cite{boylestad2015}.
It is built from two back-to-back p--n junctions sharing a thin \textit{base} region,
with \textit{emitter} and \textit{collector} regions on either side. NPN types (n-type
emitter and collector, p-type base) are more common. The BJT was invented at Bell Labs
in 1947 and was the primary amplifying device in all electronics for the following two
decades until the MOSFET gradually took over in digital applications.

\labimage{lab-01/assets/bjt_to92.jpg}%
  {Small-signal BJT in TO-92 plastic package. The three terminals---emitter,
   base, and collector---are identified by the flat face on the package
   and the manufacturer's datasheet.}{bjt_lab}

The BJT output ($I_C$--$V_{CE}$) characteristics, shown in Figure~\ref{fig:bjt_char},
display three distinct operating regions: the \textit{saturation} region (both junctions
forward-biased, used for switching), the \textit{active} region (base-emitter forward,
base-collector reverse, used for linear amplification), and the \textit{cutoff} region
(both reverse-biased, transistor off).

\theoryimage{lab-01/assets/bjt_characteristics.png}%
  {NPN BJT output characteristics ($I_C$ vs.\ $V_{CE}$) for five values
   of base current $I_B$. The saturation and active regions are indicated;
   $\beta = h_{FE} = \Delta I_C / \Delta I_B \approx 100$ in the active region.}{bjt_char}

\subsubsection*{Specifications}
Primary BJT parameters: DC current gain $h_{FE}$ (or $\beta$, ratio $I_C/I_B$, typically
50--500), collector--emitter breakdown voltage $V_{CEO}$, maximum collector current $I_C$,
maximum power dissipation $P_C$, and transition frequency $f_T$ (the frequency at which
current gain falls to unity) \cite{sedra2020}.

\subsubsection*{Purpose of Use}
In the \textit{active region} the BJT linearly amplifies current and voltage; in the
\textit{saturation/cutoff region} it operates as an electronic switch. The base--emitter
voltage $V_{BE}$ has a well-characterized temperature dependence of $-2$\,mV/$^\circ$C,
making BJTs useful as temperature sensors.

\subsubsection*{Applications}
BJTs are used in audio amplifiers, RF amplifiers, oscillators, relay and motor drivers,
TTL digital logic, temperature measurement circuits, and as the output stage of
operational amplifiers \cite{razavi2013}.
